\section{Valentía y sabiduría: Brain, Guts, Heart, Taste}

\subsection{Las tres dimensiones de Jensen Huang y la cuarta dimensión de Hu Yuanming}

El título del último capítulo de la entrevista, «Valentía y sabiduría», proviene del título de un artículo de Hu Yuanming en Zhihu. Ahí propone cuatro palabras clave para evaluar el talento: Brain, Guts, Heart, Taste; las tres primeras provienen de la respuesta de Jensen Huang en una conferencia en Stanford, y la cuarta la agregó él mismo.

\begin{importantbox}{{Brain / Guts / Heart / Taste}}
\begin{itemize}
\item \textbf{Brain} (mente): el problem solver más fuerte.
\item \textbf{Guts} (valentía): puede actuar rápido cuando la información es incompleta, puede aceptar el riesgo de fracasar, puede aprender del fracaso.
\item \textbf{Heart} (corazón): ser bondadoso desde el fondo; pero eso no significa ser la persona más nice. «A veces exigirle a tu equipo un estándar alto es, en realidad, la mayor bondad, porque si no marcas ese estándar el equipo va a fracasar.» En el fondo no puede haber intención de dañar a otros; debe ser una persona con principios y honesta.
\item \textbf{Taste} (gusto): la dimensión que Hu Yuanming agregó por su cuenta. Al contratar pregunta: «¿A qué juegos juegas? ¿Qué haces los fines de semana? ¿Qué libros lees?»; solo las personas con gustos e intereses afines pueden colaborar de cerca durante mucho tiempo.
\end{itemize}
\end{importantbox}

Sobre Taste, reveló que últimamente ha estado leyendo un libro de Lei Jun (雷军). Piensa que Lei Jun es un emprendedor del que vale mucho la pena aprender: en su momento fue un «trabajador modelo» en Zhongguancun, haciendo para Microsoft un sustituto de Office (la suite Pangu), y se perdió por completo la era de internet; más tarde resumió la teoría del cerdo volador: «en la boca del huracán, hasta los cerdos pueden volar». Esto sigue la misma lógica que su propia idea de que «elegir la dirección importa mucho más que esforzarse».

\subsection{Etapa uno y etapa dos: de quien resuelve problemas a quien los plantea}

Hu Yuanming propone un marco profundo para el desarrollo de la vida: \textbf{la etapa uno} y \textbf{la etapa dos}:

\begin{table}[H]
\centering
\caption{Diferencia esencial entre la etapa uno y la etapa dos}
\begin{tabular}{lll}
\toprule
\textbf{Dimensión} & \textbf{Etapa uno} & \textbf{Etapa dos} \\
\midrule
Origen del problema & Otros te plantean el problema & Tú mismo te planteas el problema \\
Claridad del objetivo & Hay un objetivo claro & No hay un objetivo claro \\
Alcance de la responsabilidad & Hacer bien lo que te corresponde & Pensar qué debe hacer todo el equipo \\
Habilidad central & Obedecer órdenes, competir al máximo por la eficiencia, perfeccionar habilidades concretas & Definir la dirección, enfrentar la incertidumbre \\
Criterio de éxito & Completar la tarea & Crear valor \\
\bottomrule
\end{tabular}
\end{table}

Este marco aplica a cada salto de la vida: de estudiante de licenciatura a estudiante de posgrado, de IC (colaborador individual) a Manager, de empleado de una empresa a emprendedor.

\begin{warningbox}{Cuanto más éxito tengas en la etapa uno, más puedes suffer en la etapa dos}
«Las habilidades de la etapa uno y las que exige la etapa dos son completamente distintas. En la etapa uno puede que se trate de obedecer órdenes, competir al máximo por la eficiencia, perfeccionar habilidades concretas. Pero al llegar a la etapa dos ya nadie te da órdenes, y ni siquiera se puede hablar de eficiencia, porque ni siquiera está definido el metric de la eficiencia. Mucha gente a la que le fue muy bien en la etapa uno se siente sumamente incómoda al llegar a la etapa dos: es como pelear contra un boss; terminas una barra de vida, el boss se transforma y empieza una segunda barra de vida, y si sigues peleando con los métodos de la primera etapa, mueres muy rápido.»
\end{warningbox}

Consejos concretos para estudiantes de doctorado:
\begin{itemize}
\item \textbf{Etapa uno} (los primeros 3 o 4 artículos): publicar los papers de manera honesta y disciplinada, «la etapa uno es competir al máximo en la ejecución». Obtener suficiente información y retroalimentación publicando muchos artículos.
\item \textbf{Etapa dos} (una vez que ya tienes una base): empezar a pensar en el why: «¿cómo salieron estos papers? ¿cuánta gente los cita? ¿forman un sistema coherente? ¿han trascendido su nicho?»; pasar de las publicaciones de relleno a hacer un trabajo con impact.
\end{itemize}

El criterio para juzgar si alguien puede llegar a ser un buen impact maker en la etapa dos: \textbf{ver si es capaz de reconocer que en el fondo es un pendejo: respetar los hechos, admitir el error cuando se equivoca, mejorar sin parar, y tener un buen role model, aunque ese role model esté muy lejos de él.}

\subsection{Explore vs. Exploit: visión de largo plazo y ejecución de corto plazo}

Hu Yuanming dedica cerca del 20\% de su tiempo a pensar en el futuro. Enfatiza que la vision no es algo que se hace una sola vez: «tal vez una palabra mejor sea visioning; tienes que estar pensando todo el tiempo en cómo es en realidad el futuro».

Su visión de futuro gira en torno a «AI for fun»:

\begin{knowledgebox}{AI for Fun: el producto definitivo de la era pos-AGI}
«Cuando llegue la AGI, todo el mundo va a perder su trabajo. Un mundo mejor es uno en el que la gente todavía pueda vivir feliz. Si con la tecnología de la AI generativa podemos crear para todos una immersive experience, y el tiempo limitado de cada persona se vuelve infinitamente rico gracias a la AI generativa —puedes volver al pasado, vivir vidas distintas—, creo que esto es, ultimately, una de las tecnologías más definitivas que el mundo necesita.»
\end{knowledgebox}

\subsection{La vida son diez grandes eventos}

Hu Yuanming tiene un marco singular para planear la vida: supongamos que empiezas a actuar a los 20 años y te retiras a los 70, con 50 años en medio. Si cada gran evento necesita 5 años, se pueden hacer 10 cosas. De esas, necesariamente la mitad va a fracasar: «si las 10 salen bien, eso significa que lo que hiciste tenía muy poco riesgo, y no era innovación de verdad».

Por ahora se pone una calificación de avance de «0.5 + 0.5»: haber llevado Taichi (太极) hasta donde está ahora vale 0.5, y haber llevado Meshy más allá del PMF vale otro 0.5. Sacar bien Taichi 2.0 es el siguiente 0.5, y llevar Meshy hasta la salida a bolsa es otro 0.5. Después de la salida a bolsa, ir por «el siguiente 1».

\subsection{Su manera de entender el fracaso}

\begin{importantbox}{Qué es el fracaso de verdad}
«¿Y qué si fracaso? ¿Me voy a morir? No. Mientras siga vivo, no he fracasado. Cualquier fracaso es solo algo temporal. \textbf{El fracaso de verdad} es dejar de querer innovar después de uno o dos intentos: que una serpiente te muerda una vez y durante diez años le tengas miedo hasta a una cuerda de pozo, y termines haciendo con disciplina solo cosas conservadoras. Si de 10 intentos de innovación 1 sale bien, ya es bastante bueno; fracasar es completamente normal.»
\end{importantbox}

\subsection{Consejos para los jóvenes}

\textbf{Para quienes quieren hacer un doctorado}: procurar encontrar en el camino algo que de verdad amen. «He visto a demasiada gente hacer un doctorado solo por el grado, y al final tampoco les gusta lo que hacen. Encontrar lo que amas no es algo que encuentres haciendo un doctorado: yo, antes de hacer el doctorado, ya había logrado decenas de artículos en SIGGRAPH, así que sabía que me iba a encantar.»

\textbf{Para quienes quieren emprender}:
\begin{enumerate}
\item \textbf{Get ready}: no se trata de tener claro qué vas a hacer, sino de tener la capacidad de soportar el fracaso: que tus padres estén sanos, que tu familia esté estable, poder pensar en la empresa las 24 horas.
\item \textbf{Take action}: la experiencia no importa tanto, muchas cosas se pueden aprender rápido por ensayo y error. «Si fracasas, fracasas, ¿y eso qué?»
\item Su propia mentalidad en ese momento era: «Si me gradúo del PhD en tres años y medio, aunque desperdicie dos años eso equivale a graduarme en cinco años y medio, ¿y eso qué?»
\end{enumerate}

\subsection{¿Todavía necesita un CEO escribir código?}

Hu Yuanming todavía dedica algunas horas por semana a escribir código (cerca del 10\% de su tiempo), sobre todo para ayudar al equipo a hacer debug y review del código. Citó lo que dijo Lisa Su, CEO de AMD: «su mayor fortaleza frente a un MBA es que actually understand cómo funciona la tecnología».

\begin{warningbox}{El peligro de que el CEO no escriba código}
«Si un CEO no escribe código, es muy fácil que se vuelva arrogante y crea que todo son problemas fáciles de resolver. Pero si de verdad te pones a ver cada día algunos de los "problemas de mierda" que la gente tiene que resolver, te das cuenta de que la cosa es bastante complicada. Un buen Founder siempre es hands-on: Lei Jun, Jensen Huang y Elon Musk son así. No necesariamente se trata de escribir código; puede ser diseñar el producto, hablar directamente con los proveedores, o tener un entendimiento profundo de la tecnología.»
\end{warningbox}

\subsection{Cápsula del tiempo: una palabra}

Al final de la entrevista, el conductor le pidió a Hu Yuanming que dejara una cápsula del tiempo para 2025. Él dijo una sola palabra: \textbf{valentía}.

«Espero que todos los que vean este Podcast tengan valentía.»

¿Qué haría si dejara de emprender? Dice que \textbf{escribiría un libro}: haría una «versión remasterizada en alta definición» del curso Games201, y escribiría un tutorial sobre un hybrid simulator de Simulation + Neural Network. «En su momento lo di demasiado apurado: todos los días de 8 a 9 de la mañana, y luego me iba a hacer mi práctica profesional en NVIDIA.»

\subsection{La distribución global de Meshy y la contratación de personal}

Meshy es una empresa globally distributed: el propio Hu Yuanming está en Silicon Valley, y el equipo está repartido entre Beijing, Shanghái, Shenzhen y otras partes del mundo. La empresa está pasando poco a poco del modelo work-from-home a un esquema híbrido: al menos 4 días a la semana en la office (en la office de China, los miércoles es teletrabajo), porque «discutir en persona es más eficiente».

En cuanto a la contratación, Meshy actualmente se enfoca en contratar AI Researcher, Machine Learning Engineer y Compiler Engineer (como reserva para Taichi 2.0). A los becarios se les paga con un salario global unificado (global pay), con un máximo de 1,000,000 de yuanes o 140,000 dólares al año, y el propio CEO participa directamente en guiar a los becarios en producto e investigación.

Esta manera de «que el CEO mismo guíe a los becarios» refleja la insistencia de Hu Yuanming en ser hands-on: él considera que esto es la ventaja competitiva central de una empresa con ADN técnico: \textbf{si el CEO no entiende los detalles técnicos, ¿por qué habría de trabajar para ti el resto de la gente?}

\subsection{Resumen de la sección}

Esta sección concentra el núcleo de la filosofía de vida de Hu Yuanming: la visión de talento en cuatro dimensiones Brain/Guts/Heart/Taste va más allá del marco tradicional de «inteligencia + esfuerzo»; el modelo de salto de la etapa uno a la etapa dos explica por qué muchos colaboradores individuales excelentes no logran ser buenos Managers; la redefinición del fracaso («mientras siga vivo, no he fracasado») da una base psicológica para la innovación constante. Y la palabra que le deja al mundo —valentía— es tanto una síntesis de su propio camino como una expectativa para quienes vienen después.
